\begin{abstract}
We characterize the exact regular local asymptotic minimax squared-error bound for estimating the average treatment effect in an independently sampled randomized trial with binary outcomes, known interior assignment probability, interior arm means, and a fixed positive local privacy budget. Each subject releases one privatized assignment--outcome record, and release rules may adapt to earlier outputs on arbitrary measurable spaces. The optimal stationary contrast variance is the reciprocal of a finite, fourteen-pattern staircase saddle that profiles out the common-mean nuisance direction, and it lower-bounds local asymptotic minimax squared-error risk for every sequential private procedure. The attainment results assume a fixed measurable strong-saddle selector satisfying the stated conditions. Fixing any such selector, a private-pilot affine estimator attains the bound within the regular class and supplies consistent variance estimation and pointwise asymptotically valid Wald intervals, including at optimal-support changes. An attaining stationary release always exists with at most five active outputs; certified open neighborhoods have exact minimum cardinalities of five and three. Balanced assignment and complementary arm means yield an optimal binary randomized-response release. A comparison with the published known-assignment custom-Laplace sample mean translates the variance improvement at a specified balanced design point into prospective interval-length and recruitment calculations.
\end{abstract}

\section{Introduction}

This paper identifies the exact regular local asymptotic minimax squared-error constant for estimating a treatment effect when both treatment assignment and outcome are locally private. The experiment consists of independently sampled subjects in a randomized binary trial, with a known assignment probability strictly between zero and one, interior outcome probabilities in both arms, and a fixed positive privacy budget. Each subject contributes one private assignment--outcome record. Release rules may depend on preceding private outputs and may take values in arbitrary measurable spaces. Within this model, the stationary optimal contrast variance is the reciprocal of a finite nuisance-profiled staircase optimization. Fixing any measurable strong-saddle selector satisfying the stated conditions on the interior parameter domain, a private-pilot estimator attains the constant and supports pointwise Wald inference.

Randomization identifies the average treatment effect as the difference between the treatment-arm and control-arm outcome probabilities. Joint local privatization determines how much information about that difference reaches the analyst. A release must satisfy the privacy inequality across all four assignment--outcome pairs, including pairs that differ in both coordinates. The resulting design problem concerns the scalar treatment-effect contrast within a two-parameter model: both arm means vary, and their common shift is a nuisance direction. Efficient release selection therefore depends on how the channel preserves information about the contrast after accounting for that nuisance variation.

Our first result connects this statistical objective to finite channel geometry. Under the stated interiority and fixed-privacy conditions, \cref{obj:thm:finite-oracle} characterizes the smallest stationary contrast variance over all measurable output spaces. Its reciprocal is obtained by maximizing nuisance-profiled information over weights on fourteen nonconstant staircase patterns. For each fixed profiling direction, information is linear in these weights; minimization over the common-mean direction produces the saddle objective. Equivalently, the optimized information is the minimum of a finite upper envelope of convex quadratics. This representation gives an exact finite optimization at each specified pair of arm means, assignment probability, and privacy budget.

The stationary optimum also governs the sequential experiment. Conditional on a fixed measurable selector satisfying the strong-saddle conditions in \cref{obj:def:pilot-estimator}, \cref{obj:thm:sequential-minimax} establishes a local asymptotic squared-error lower bound for every sequence of sequential private protocols and transcript-measurable estimators. The criterion takes the worst risk over bounded admissible local perturbations, then the lower asymptotic limit, and finally the supremum over perturbation radii. The same constant is attained within the regular class by a private-pilot construction with pilot size equal to the integer part of the square root of the total sample size. Thus history adaptation and arbitrary measurable releases are incorporated into the exact efficiency comparison.

The construction first estimates the arm means using four-category randomized response and then selects a staircase channel for the remaining subjects. Its treatment-effect estimate combines the pilot contrast with an affine correction computed from the main-sample releases. Under a fixed measurable strong-saddle selector, \cref{obj:thm:pilot-attainment} establishes exact conditional centering given the private pilot history, regular asymptotic normality, consistent empirical variance estimation, and pointwise asymptotic Wald coverage. These conclusions hold for every pilot schedule satisfying \cref{obj:ass:pilot-diverges,obj:ass:pilot-sublinear}. They apply throughout the interior parameter domain, including points where the optimal channel support changes.

The optimal output alphabet has a corresponding statistical structure. \Cref{obj:thm:five-output-upper-bound} guarantees an attaining stationary staircase release with at most five active outputs at every interior design point. At assignment probability \(3/10\) and privacy budget \(\log 3\), \cref{obj:thm:support-transition} identifies open neighborhoods of arm means \((1/2,1/2)\) and \((1/10,1/10)\) in which the minimum attaining cardinalities among all stationary channels are five and three, respectively. Under balanced assignment and complementary arm means, \cref{obj:thm:balanced-reduction} instead yields an optimal two-output randomized-response release. These results establish specific optimal release cardinalities under contrast-specific efficiency at selected outcome prevalences and designs.

The analysis connects two established strands of private inference. \Citet[Theorems 3.3 and 3.5, Remark 3.6, and Section 4]{Steinberger2024} develops sequential asymptotic efficiency theory under his stated measurable-space and information-limit conditions, and identifies scalar functionals of inverse information as objectives in multiparameter models. Here the four-record trial model permits an exact nuisance-profiled contrast calculation over arbitrary measurable sequential output spaces. The staircase geometry follows the extremal-mechanism framework of \citet[Theorems 2 and 4]{KairouzOhViswanath2016}; nuisance profiling gives the present optimization its concave weight objective and associated support structure. The exact fixed-privacy constant complements the statistical-rate analyses of \citet{DuchiJordanWainwright2018,Rohde2020,Duchi2024}.

For a direct causal benchmark, we evaluate the known-assignment custom-Laplace release and sample-mean estimator of \citet[Sections 3.1 and 3.3, Theorem 7, and Appendix A.5]{OhnishiAwan2025}. At assignment probability \(1/2\), both arm means \(1/2\), and privacy budget one, \cref{obj:thm:laplace-comparison} establishes a strict variance improvement for the efficient procedure. The implied asymptotic half-length ratio is approximately \(2.69\), with the custom-Laplace sample mean in the numerator; the corresponding sample-size ratio at equal precision is approximately \(7.26\). We illustrate these quantities through a prospective binary endpoint defined by meeting the \(80\%\) attendance threshold of the Colombian cash-transfer program studied by \citet{BarreraOsorioEtAl2011}.

The paper next reviews related work, defines the private trial model, derives the efficiency bound and finite optimization, develops private-pilot inference, studies optimal release cardinality, and gives the interval comparison and prospective design calculation. Discussion is followed by the proof exposition in \cref{sec:proofs} and the appended proofs.
